% Appendix: Proofs for Section 5 (drafted 2026-11-23, Phase 4 unit 4).
% Holds the full NCL finite-step argument (OBJECTIONS item 3: both
% directions, every feasible step below the cutoff, several or negative
% summands, the common represented orientation, mixtures legal,
% preload/toggle/unload, loops counting twice), the control witness of
% Remark 5.7 (rem:control), the algebraic and real NCL transfers, the rational-root
% compiler and the fixed quartic layout. The Heintz-Roy-Solerno upper
% bound is the last subsubsection (moved here from appendix-b-proofs.tex).
% Sources: lean/SingleAdditiveLicence.md sections 1-4 and 7;
% lean/BoundaryCorollaries.md sections 1, 2 and 4; lean/BoundedProbes.md
% sections 1 and 3; lean/ComplexityBoundary.md section 3.
% Lean (table only): ConstraintLogic legal_of_represents,
% flips_of_represents, common_orientation, represents_corner,
% Graph.flip_segments_lawful; SingleAdditiveLicence ncl_chain_iff,
% ncl_fine_iff, ncl_endpoints_distinct, control_replacement, Control.fine,
% Control.verdict_changes; AlgebraicConstraintLogic chain_of_flips,
% flips_of_chain, fine_iff_flips; RealAlgebraicBoundary
% real_path_iff_flips, flips_of_real_path; ResponseBudgets (lift_chain,
% decode_chain, fine_iff_box, no_native_chain_of_selector, square_rows);
% QuarticSelector (given a RowRealization).

\subsection{Proofs for Section~\ref{sec:hard}}\label{app:hard-proofs}\label{app:ncl-proofs}

% ---------------------------------------------------------------------------
\subsubsection{Size of the construction}\label{app:ncl-size}

% compressed 2026-11-24: (a) the label count and one-value register, given
% in Section 5.2, are not re-listed; (b) the size bound is computed in one
% sentence from the listed counts (entry count, denominator, numerator
% length).
The construction has $|I|=|E|+1$ probes,
$|E|+2\le|X|<2(|E|+2)$, $2|E|+1$ word tests and $|I|+|E||I|+|V|$ value
tests. As $2^d-1\ge|I|$, listing $X$ and its first $|I|$ nonzero points in
binary order gives the~$a_i$, and $\chi_i(g)=(-1)^{a_i\cdot g}$ takes
$O(d|I||X|)$ Boolean operations. For
distinct $a_i,a_{i'}$ the nonzero character $\chi_i\chi_{i'}$ pairs points
of opposite sign, so the probes are balanced and orthogonal. At both
endpoints the eigenvalues have denominators dividing $8|I|$ and the label
probabilities are $\delta$ or~$2\delta$. Over the common denominator
$8|I||X|(|M|+1)$, each of the $|M||X|^2$ entries
$\theta_m(1+\sum_i\alpha_{mi}\chi_i(g)\chi_i(h))/|X|$ therefore has a
numerator summing $O(|I|)$ integers of binary length
$O(\log(2+|E|+|V|))$. With
$1/\varepsilon_0=(|M|+1)\cdot4|I|\cdot(2|X|+1)$, the output has binary
length $O((|E|+|V|+1)(|E|+1)^2\log(2+|E|+|V|))$ and is computed in
polynomial time.

% ---------------------------------------------------------------------------
\subsubsection{The finite-step argument}\label{app:ncl-transfer}

% compressed 2026-11-24: (e) sentences merged in the set-up, (M), (J), (L)
% and both lemma proofs. Every element of OBJECTIONS item 3 is kept: both
% directions, every feasible step below epsilon_0, several or negative
% summands, the common represented orientation, mixtures legal,
% preload/toggle/unload, and loops.
Configurations here are feasible at the profile~$p$ of the construction,
so Lemma~\ref{lem:feasible} applies. For such $A$ and $e\in E$, let $C_A(e)\subseteq\{0,1\}$ contain $0$ if $x_{A,e}\neq0$ and $1$
if $y_{A,e}\neq0$; it is nonempty by~\eqref{eq:load}. Thus $A$ represents
$\sigma$ exactly when $\sigma(e)\in C_A(e)$ for every~$e$, so its
represented orientations are all mixtures, one choice per edge, of the
sets~$C_A(e)$.

\emph{(M) Every represented orientation is legal.} Let $A$ represent
$\sigma$, and fix~$v$. Each edge $e$ with tail $v$ under $\sigma$ has the
end $v_e^{\sigma(e)}=v$ available, so its selector equation
in~\eqref{eq:selector} gives $\alpha_{\mathsf t_ve}=\wt(e)/8$. Distinct
edges give distinct coordinates, and a loop at $v$ has one coordinate and
one tail, although it contributes twice to~$W(v)$. Summing,
\eqref{eq:vertex-budget} gives the tails at~$v$ weight at most $W(v)-2$,
so the heads carry at least two. This bound uses neither the signs of the
summands nor the other coordinates of~$\mathsf t_v$, which may carry
unused charge.

\emph{(J) Represented orientations are joined inside one configuration.}
If $A$ represents $\rho$ and $\rho'$, change the edges where they differ one
at a time. Each intermediate orientation agrees at every edge with $\rho$
or~$\rho'$, so $A$ represents it and it is legal by~(M), and consecutive
ones differ by one flip.

\emph{(L) Lifts.} At $A_\sigma$ one summand vanishes at each edge, so
$A_\sigma$ represents only~$\sigma$. For legal $\sigma$, $A_\sigma$ meets
Lemma~\ref{lem:feasible}. Each tail's vertex label carries the charge and
the other summand is zero, so the selector equations hold. The vertex
label at $v$ uses its tail weight divided by~$8$ of the edge-probe budget,
whose limit is $(W(v)-2)/8$. The
summand labels spend at most $|I|k=1/4$, the load label $1/8$, and each
selector at most~$1/4$. So distinct legal orientations give distinct
feasible endpoints.

\begin{proof}[Proof of Lemma~\ref{lem:decode}]
Let $e\in E$. If $C_A(e)\cap C_B(e)=\emptyset$, the two sets are opposite
singletons, say $\{0\}$ and~$\{1\}$; then $y_{A,e}=0$, so $x_{A,e}=k$
by~\eqref{eq:load}, while $x_{B,e}=0$, contradicting
$|x_{A,e}-x_{B,e}|<k$. Choosing $\rho(e)\in C_A(e)\cap C_B(e)$ at each edge
independently gives an orientation $\rho$ represented by both, whatever
the other coordinates do, whether one or many change, and whatever the
signs of the summands.

Now let $A_\sigma=B_0,B_1,\dots,B_n=A_\tau$ satisfy the bound at each step.
If $n=0$, then $A_\sigma=A_\tau$ and $\sigma=\tau$ by~(L). Otherwise choose
$\rho_s$ represented by $B_{s-1}$ and $B_s$ $(1\le s\le n)$, and put
$\rho_0=\sigma$ and $\rho_{n+1}=\tau$; by~(L), $\rho_1=\sigma$ and
$\rho_n=\tau$. Each $B_s$ $(0\le s\le n)$ represents $\rho_s$ and
$\rho_{s+1}$, so (J) joins them by legal flips, possibly several within
the one configuration~$B_s$, and concatenation gives legal flips from
$\sigma$ to~$\tau$.
\end{proof}

% compressed 2026-11-24: (e) sentences merged; the vertex-budget argument
% is stated as one sentence tracking each total. The preload, toggle and
% unload segments, the loop case, every feasibility check and the explicit
% subdivision bound are kept.
\begin{proof}[Proof of Lemma~\ref{lem:lift}]
Let $f$ be the flipped edge, with $\sigma(f)=0$ (the other case is
symmetric), so $a=v_f^0$ is the old tail and $b=v_f^1$ the new one. The
lifts $A=A_\sigma$ and $A'=A_\tau$ differ only in $(x_f,y_f)$, which is
$(k,0)$ at $A$ and $(0,k)$ at~$A'$, and, if $a\neq b$, in
$(\alpha_{\mathsf t_af},\alpha_{\mathsf t_bf})$, which is $(\wt(f)/8,0)$ at
$A$ and $(0,\wt(f)/8)$ at~$A'$. Let $A^+$ be $A$ with
$\alpha_{\mathsf t_bf}=\wt(f)/8$, and $A'^+$ be $A'$ with
$\alpha_{\mathsf t_af}=\wt(f)/8$. At fixed label probabilities and
register, move affinely in the eigenvalues from $A$ to $A^+$ (preload),
then to $A'^+$ (toggle, moving $(x_f,y_f)$ from $(k,0)$ to $(0,k)$), then
to $A'$ (unload). For a loop, $a=b$, the charge stays $\wt(f)/8$, and the
preload and unload are constant.

Every point meets Lemma~\ref{lem:feasible}. The moved coordinates are not
calibrated. The load equation holds, with $x_f,y_f\in[0,k]$ on the toggle,
so the budgets of $\mathsf x$ and $\mathsf y$ stay at most $|I|k$. The
selector equations at~$f$ hold: on the preload, $y_f=0$ and the charge at
$a$ is present; on the toggle, both charges are; and on the unload, $x_f=0$
and the charge at $b$ is present. Other edges are unchanged, so only the
vertex budgets at $a$ and $b$ change. The total at $b$ rises along the
preload from its value at $A$ to that at~$A'$; the total at $a$ falls
along the unload from its value at $A$ to that at~$A'$; and the toggle
fixes all charges. So each total stays between its two endpoint totals,
which satisfy~\eqref{eq:vertex-budget} by~(L).

Every entry $P_m((g,r),(h,r))=\theta_mQ_m(g,h)$ is affine in the
eigenvalues. On each segment every label changes at most one eigenvalue,
by at most $1/4$, so dividing each segment into
$n\ge1/\varepsilon$ equal parts gives rational feasible configurations
at~$p$ whose consecutive entries differ by at most $1/(4n)<\varepsilon$.
\end{proof}

% compressed 2026-11-24: (c) the eigenvalue-change estimate, which repeated
% the signed row sum of Lemma transfer(a) (sum_h chi_i(h) P_m = theta_m
% alpha_mi chi_i(g), |X| terms, divide by delta), now cites Lemma
% transfer(a) with N=|X|; since 1/epsilon_0 has the factor 2|X|+1, that
% bound is already below k. (e) merges. The three-way equivalence is kept.
\begin{proof}[Proof of Lemma~\ref{lem:cutoff}]
First, a step below $\varepsilon\le\varepsilon_0$ changes every eigenvalue
by less than~$k$: with the one-value register, $N=|X|$ in
Lemma~\ref{lem:transfer}(a), which bounds the change by
$2|X|\varepsilon_0/\delta=2|X|k/(2|X|+1)<k$.

Along legal flips $\sigma=\rho_0,\dots,\rho_n=\tau$, concatenating the
sequences of Lemma~\ref{lem:lift} (empty if $n=0$) gives fine reachability
of $A_\tau$ from~$A_\sigma$, hence a sequence at~$\varepsilon$. Conversely,
a sequence at $\varepsilon\le\varepsilon_0$ is one at $\varepsilon_0$, all
at~$p$ (Remark~\ref{rem:inputs}(ii)); by the first paragraph its
consecutive summands differ by less than~$k$, and Lemma~\ref{lem:decode}
gives legal flips from $\sigma$ to~$\tau$.
\end{proof}

% ---------------------------------------------------------------------------
\subsubsection{The control of Remark~\ref{rem:control}}\label{app:ncl-control}

% compressed 2026-11-24: (c) the closing subdivision cites rational
% subdivision; (d) the closing sentence "In the original instance, Lemma
% cutoff gives the verdict no ..." is cut, as Remark rem:control states it;
% (d) the K_{3,3} paragraph no longer repeats the remark's AND-vertex
% count ("each vertex meets one matching edge of weight two and two edges
% of weight one"); it keeps the inflow computation, the no-legal-flip check
% and distinctness.
The replacement test $\mathsf{uu}\sim\mathsf{uu}$ passes for every parameter
vector, and all other inputs, held tests and admissible sets are
unchanged, so $A_\sigma$ and $A_\tau$ stay feasible at the same profile.
Fix their label probabilities, $2\delta$ at $\mathsf u$ and $\delta$
elsewhere, and the register. For rational $t\in[0,1]$, let $A_\sigma(t)$
be $A_\sigma$ with $x_i,y_i$ multiplied by $t$ at every probe, the reserve
included. It is feasible in the replacement: the selector tests keep equal
word probabilities, their probe equations scale by~$t$, the summand
budgets decrease, and no calibration involves $\mathsf x$ or~$\mathsf y$.
At $t=0$ every selector equation reads $0=0$, so the
$\alpha_{\mathsf t_ve}$ $(e\in E)$ move affinely from their values at
$A_\sigma$ to those at $A_\tau$ within the convex budgets, with
$\alpha_{\mathsf t_v\ast}$ fixed by its calibration. Finally, follow
$A_\tau(t)$ from $t=0$ to~$1$. These segments are affine in the
eigenvalues at fixed label probabilities and register, so rational
subdivision gives sequences at every positive
tolerance.


% ---------------------------------------------------------------------------
\subsubsection{The algebraic and real transfers of Section~\ref{sec:real-ncl}}\label{app:ncl-algebraic}

% compressed 2026-11-24: (c) the decoding over R (common orientation for
% steps below one, then (M), (J), (L) as in the proof of Lemma decode) is
% stated once in the first paragraph and cited by both converses, rational
% and real, instead of being restated; (e) sentences merged. Both
% directions of both transfers are kept.
In~\eqref{eq:ncl-system}, $|x_{e0}|,|x_{e1}|\le1$ and $x_{e0}+x_{e1}=1$
force $x_{eb}\in[0,1]$, and $|z_{ve}|\le W(v)-2\le4$, so the real, hence
the rational, feasible set is bounded. Side $b$ of $e$ is available at a
feasible point if $x_{eb}\neq0$; every edge has an available end, and an
orientation is represented if every tail is available. (M), (J) and (L)
hold with the same proofs over~$\RR$, with the charge $\wt(e)$ for
$\wt(e)/8$ and the budget $W(v)-2$ for $(W(v)-2)/8$; so the endpoint of a
legal orientation is feasible and represents only it. Two feasible points
with $|x_{e0}-x'_{e0}|<1$ at every edge represent a common orientation, by
the proof of Lemma~\ref{lem:decode} with $k=1$. So, as there, real feasible
points from the endpoint of $\sigma$ to that of~$\tau$ with consecutive
sup distances below one give legal flips from $\sigma$ to~$\tau$.

Take a flip from side~$0$ to side~$1$, with old and new tails $a$
and~$b$. If $a\ne b$, raise $z_{bf}$ to $\wt(f)$, move the summands from $(1,0)$ to
$(0,1)$, and lower $z_{af}$ to~$0$; if $a=b$, keep $z_{af}=\wt(f)$ and move
only the summands. Reverse these segments for a flip from side~$1$ to
side~$0$. The checks are those of the proof of Lemma~\ref{lem:lift}, with
unscaled charges and unit summands, and the coordinates are affine in the
parameter, so rational subdivision gives feasible rational sequences with
sup-norm steps below any $\varepsilon>0$. Conversely, fine connectivity gives such a sequence with
steps below~$1$, hence legal flips. So the endpoints are finely connected
exactly when legal flips lead from $\sigma$ to~$\tau$, and the promises of
Corollary~\ref{cor:algebraic} hold. The system has $2|E|+|V||E|$
coordinates, $2|E|$ binomial tests, $|E|$ normalized affine trinomials and
$2|E|+|V|$ budgets. The system and endpoints are written in polynomial
time from the weights, capacities $W(v)-2$, and $0$ and~$1$.

For real path components: the tests and budgets are continuous, so the
real feasible set is closed, and each segment, feasible at rational
parameters, is feasible on its whole real interval; so legal flips give a
piecewise-linear real path. Conversely, a continuous
path $\gamma:[0,1]\to\FR$ joining the endpoints is uniformly continuous,
so for some $n$ the samples $\gamma(s/n)$, $0\le s\le n$, have consecutive
sup distances strictly below one, and give legal flips as above. Hence
real path-component membership on these bounded instances is
PSPACE-hard.

% ---------------------------------------------------------------------------
\subsubsection{Proof of Theorem~\ref{thm:rational-roots}}\label{app:rational-roots}

% compressed 2026-11-24: (e) sentences merged in all four parts; (c) the
% transfer cites rational subdivision. Every gadget, inequality, label,
% calibration, budget check and size count is kept.
\emph{The gated system.} Build each coefficient of $f$ from its binary
digits by addition and doubling from $0$ and~$1$, negating for its sign,
raise each variable to its binary exponents by repeated squaring, and
multiply and sum. This gives $O(\mathcal L+1)$ gates, to which every input
assignment extends uniquely, the last evaluating~$f$. With the
homogenized gate equations of Section~\ref{sec:rational-roots}, a gated
residual is $s(2-s)$ times the residual of a gate equation, of the output
equation $z_\gamma=0$, or of $\xi^2-2-\zeta^2\psi$. At a rational feasible
point with $0<s<2$, the switch factor is nonzero, so the ungated residuals
vanish. If $\zeta=0$, the slack equation would give $\xi^2=2$, which has
no rational solution. Thus $\zeta>0$. Dehomogenizing the gate and output
equations shows that $v_\gamma=z_\gamma/\zeta$ satisfies the original
gate equations with output zero. Conversely, extend a rational root $v$ to the gates, put
$\zeta=(1+\sum_\gamma|v_\gamma|)^{-1}$ and
$z_\gamma=\zeta v_\gamma\in[-1,1]$, choose rational $\xi$ strictly between
$\sqrt2$ and $\sqrt2+\zeta^2/4$, and put $\psi=(\xi^2-2)/\zeta^2$. Then
$1<\xi<2$ and, as $\xi+\sqrt2<4$, $0<\xi^2-2<\zeta^2$, so $0<\psi<1$, and
these values satisfy every equation for all~$s$. Every residual vanishes on
the faces $s\in\{0,2\}$, so they are full boxes, and affine segments
inside them join the supplied endpoints to the ends of the segment
$s\in[0,2]$. Without a rational root, every rational feasible
point lies on a face.

\emph{Inequalities.} Normalize affinely: $z_\gamma=2t-1$, $\zeta=2t$,
$s=2t$, $\xi=2t+1$ and $\psi=2t-1$. Expand each residual, of bounded
degree with boundedly many terms per gate, giving each partial product its
own coordinate in $[0,1]$. Impose each product $t''=tt'$ by
$\bar t=(t+t')/2$, the squares $t^2$, $t'^2$, $\bar t^2$ and the affine
equality $2t''=4\bar t^2-t^2-t'^2$, which determine $t''$ uniquely. This
leaves affine equalities and squares $t'=t^2$ in $O(\mathcal L+1)$ coordinates of
$[0,1]$, of bounded arity with bounded-length coefficients, with unique
polynomial lifts on solutions and the switch coordinate $t_s=s/2$
unchanged. Write an affine equality as $\sum_jc_jv_j=c_0$ with $c_j\ge0$
and each $v_j$ either $t_j$ or $1-t_j$. Impose $\sum_jc_jv_j\le c_0$
and $\sum_jc_j(1-v_j)\le\sum_jc_j-c_0$; both bounds are nonnegative because
the lifted endpoints below satisfy the equality. Impose $t'=t^2$ by
$t^2+(1-t')\le1$ and $(1-t)+t(1-t)+t'\le1$, that is, $t^2\le t'$ and
$t'\le t^2$, as $(1-t)+t(1-t)=1-t^2$. Read repeated occurrences of a
coordinate from copies $t_i$ with $t_i+(1-t_j)\le1$ and
$t_j+(1-t_i)\le1$, and add $|t_j|\le1$ and $|1-t_j|\le1$ for each~$j$.
Every inequality is then
$\sum_jc_{\rho j}|t_j^{a_j}(1-t_j)^{b_j}|\le c_\rho$ with
$c_{\rho j},c_\rho\ge0$, $a_j+b_j\le2$ and at most one term per coordinate.
The endpoints have $t_s=0$ and $t_s=1$, respectively, and every other
normalized source coordinate zero. Their gadget coordinates follow in the
fixed order of partial products, squares and copies; the construction below
gives their label probabilities and eigenvalues.

\emph{The model.} Let $J$ index the coordinates and copies. Take probes
$I=J\sqcup\{\ast\}$ on the binary base of Section~\ref{sec:construction}
for $|I|$ probes, a one-value register, only level~$0$, identity reversal,
$\kappa=1/2$, $H=4$, $\delta=1/(|M|+1)$ and $k=1/(4|I|)$. The labels are
$\mathsf u,\mathsf x,\mathsf y$, a pad $\mathsf p$, an indicator
$\mathsf o_j$ for each $j\in J$, and for each inequality $\rho$ a label
$\mathsf r_\rho$ and coefficient labels $\mathsf c_{\rho j}$ $(j\in J)$.
Put $\omega_\rho=1/(2(1+c_\rho+\sum_jc_{\rho j}))$, with $c_{\rho j}=0$ and
$a_j=b_j=0$ for $j$ outside the inequality. The calibrations are $u_i=k/2$
and $\alpha_{\mathsf pi}=k$ at every probe; $x_\ast=k$;
$\alpha_{\mathsf o_ji}=k$ if $i=j$ and $0$ otherwise;
$\alpha_{\mathsf c_{\rho j}i}=\omega_\rho c_{\rho j}$ if $i=j$ and $0$
otherwise; and $\alpha_{\mathsf r_\rho\ast}=1/2-\omega_\rho c_\rho$. The
word tests are $\mathsf u\sim\{\mathsf x,\mathsf y\}$ and, for every
inequality~$\rho$ and $j\in J$, the test of
Section~\ref{sec:rational-roots}, of length $2+a_j+b_j\le4$. The only
admissible set is the full list, and the endpoints are the lifts of the two
corners, at the common full-support profile with the unique level-zero
partition and every test held.

As in Lemma~\ref{lem:engine} and Proposition~\ref{prop:length-one}, a
feasible configuration has $\theta_{\mathsf u}=2\delta$, every other label
probability~$\delta$ and $x_i+y_i=k$; so $y_\ast=0$, and with $t_j=x_j/k$,
$y_j=k(1-t_j)$. The single-word tests avoid $\mathsf u$, so their words
have equal probabilities; at probes $i\neq j$ both sides vanish, and at
$j$ cancelling $k$ gives
$\alpha_{\mathsf r_\rho j}=\omega_\rho c_{\rho j}t_j^{a_j}(1-t_j)^{b_j}$. As
$\omega_\rho c_\rho<1/2$, the budget of $\mathsf r_\rho$ is exactly
inequality~$\rho$. The other budgets hold whenever $t\in[0,1]^J$:
$\mathsf x$ and $\mathsf y$ spend at most $|I|k=1/4$, $\mathsf p$ also
$1/4$, $\mathsf u$ spends $1/8$, and $\mathsf o_j$ and $\mathsf c_{\rho j}$
at most $1/2$. So the feasible configurations at the common profile
correspond exactly to the solutions $t\in[0,1]^J$ of the inequalities,
through eigenvalues polynomial in~$t$.

\emph{Transfers and size.} If $f$ has a rational root, the three segments
lift through the gadgets and the model to configurations with eigenvalues
polynomial in the segment parameter and rational at rational parameters.
The tables are affine in the eigenvalues, so rational subdivision gives
sequences at every positive tolerance. If
$f$ has none, every feasible rational configuration has $t_s\in\{0,1\}$,
while the endpoints have $t_s=0$ and $t_s=1$. A sequence at
$\varepsilon\le\varepsilon_0=\delta k/(2|X|+1)$ would contain a step
changing $x$ at the probe of $t_s$ by~$k$, which Lemma~\ref{lem:transfer}(a) excludes, since the register has one value and $2|X|\varepsilon_0/\delta<k$. There are $O(\mathcal L+1)$ coordinates,
copies and inequalities, and one word test and coefficient label per
inequality and coordinate, so $O((\mathcal L+1)^2)$ of each and $O((\mathcal L+1)^3)$
calibration entries. All constants and endpoint coordinates have
polynomial binary length, $|X|<2(|I|+1)$, and
$1/\varepsilon_0=(|M|+1)\cdot4|I|\cdot(2|X|+1)$, so the map runs in
polynomial time and never computes a root.

% ---------------------------------------------------------------------------
\subsubsection{Proof of Corollary~\ref{cor:quartic}}\label{app:quartic}

% compressed 2026-11-24: (a) the two gated equations are cited from the
% proof idea of Corollary quartic instead of being restated; (e) merges.
% The slot layout, the counts, the corner lift, the calibration length,
% padding, slopes and the conditional composition are kept.
Let $f(v)=\sum_{|a|\le4}c_av^a$ in $N$ variables, listing every exponent
vector of total degree at most four, with zero coefficients kept. In the
box $z\in[-1,1]^N$, $\zeta\in[0,2]$, $\xi\in[1,3]$, $\psi\in[-1,1]$,
$s\in[0,1]$, take the two gated equations of the proof idea of
Corollary~\ref{cor:quartic}, and the endpoints with $s=0$ and $s=1$ and
every other normalized coordinate zero. Both faces are feasible. At a
rational point with $0<s<1$, $\zeta\neq0$ as before, and
$F(\zeta v,\zeta)=\zeta^4f(v)$ gives the root $v=z/\zeta$; conversely a
root gives the three segments of Appendix~\ref{app:rational-roots}, with
$z=\zeta v$.

After normalization both residuals have degree at most six in the $N+4$
coordinates. Compute the monomial slots of degrees one to six in these
coordinates by consecutive products in a fixed order, each through the
gadget of Appendix~\ref{app:rational-roots}. Two final affine
equalities set the constant term plus the weighted sum of the slots to
zero, with weights determined by the $c_a$; a zero coefficient keeps its
slot. Convert to inequalities as before, with a distinct copy for every
term occurrence, including those with zero coefficient, and with the two
copies of $t$ in the second square inequality. With $U$ products, $H'$
nonconstant slots and $C=29U+4H'$ term occurrences, there are $N+4+5U+C$
coordinates and copies and $2(5U+2)+2C$ inequalities besides the box
inequalities, all depending only on~$N$. So the numbers of probes, labels,
base states, word tests and value tests of the model of
Appendix~\ref{app:rational-roots} depend only on~$N$, as do $\delta$, $k$,
$|X|$ and $\varepsilon_0$. Its single-word tests have length
$2+a_j+b_j\le4$, the pure and mixed square terms having degree two.

The polynomial assignment lifts both face corners. The layout contains the
product $s\cdot s$, whose gadget values at $s=1$ are all one. Both corners
satisfy the two gated residuals, and the assignment satisfies every
product, square and copy constraint, so the compiler produces two feasible
configurations at the common profile. Each calibration takes a number of
rational operations on the coefficients bounded in terms of~$N$, so it has
binary length $O_N(\mathcal L+1)$ for input length~$\mathcal L$; a sparse input is
padded to the dense layout in polynomial time. The transfers are
those of Appendix~\ref{app:rational-roots}, at the same endpoints and the
fixed tolerance~$\varepsilon_0$. The slopes of the lift are bounded on
each segment, as every product of unit factors varies by at most its
degree times the parameter change. Composing with a fixed~$N$ from
the Bhargava--Burungale quartic theorem~\cite[Corollary~1.2(3)]{bhargavaburungale}, if it
holds, gives the conditional undecidability.

% ---------------------------------------------------------------------------
\subsubsection{The upper bound of Corollary~\ref{cor:real-complexity}}\label{app:hrs}

% compressed 2026-11-24: (e) sentences merged in all four paragraphs; the
% space bound, the degree-two translation, the unique lifts, the singleton
% endpoints, the parity count (also for coinciding endpoints) and the
% path-component step are kept.
\begin{proof}[Proof of the PSPACE upper bound in Corollary~\ref{cor:real-complexity}]
The algorithm of Heintz, Roy and Solern\'o~\cite[Section~1.1, Main
Theorem]{heintzroysolerno1994components} takes a quantifier-free formula
with dense integer coefficient vectors and outputs one formula per
semialgebraically connected component. It runs in deterministic Turing space
$(n\log D\,\sigma)^{O(1)}$ for $n$ variables, binary coefficients of
length at most~$\sigma$, and $D$ the sum of the total degrees of the input
polynomials. The input is dense, so writing out the original polynomials
can take exponential size; we first translate in polynomial time to an
explicit input of degree two.

\emph{Degree two.} Compute each monomial by repeated squaring and
multiplication (using~$1$ for the empty product), with an auxiliary
variable and equation $w=uv$ per gate. The gate count is linear in the
exponents' binary lengths and the number of factors. Each binomial test
compares its two monomial outputs by its coefficient~$c$. Auxiliary
variables $u_j\ge0$ with $u_j^2=x_j^2$ give $u_j=|x_j|$ uniquely, so each
budget becomes linear. Keep the equations $x_j+y_j=1$ and clear rational
denominators in each atom. The numbers of variables and atoms and the
binary lengths of the coefficients are polynomial, so even dense
coefficient vectors of total degree two have polynomial size.

\emph{Lifts and endpoints.} Every original feasible point has exactly one
lift, with auxiliary coordinates continuous in the original ones, so
projection is a homeomorphism from the lifted locus $G$ onto the original
real feasible set, matching path components. Conjoining
the formula of $G$ with $x=a$ specifies the lift of an endpoint $a$ as a
singleton, without writing out its coordinates, which can have
exponentially many binary digits.

\emph{Parity.} For the two singleton lifts $A$ and~$B$, apply the
component algorithm to
\[
  (G\times\{0,1\})\ \cup\ \bigl((A\cup B)\times[0,1]\bigr),
\]
which has a polynomial-size formula from the same descriptions. If $G$ has
$c$ components, the vertical intervals merge one pair of copies when the
endpoints share a component and two pairs otherwise. So the output count
is $2c-1$ in the yes case and $2c-2$ in the no case, also when the endpoints
coincide. Its parity is tallied in polynomial space while the component
descriptions are streamed. Over~$\RR$ each connected component of a
semialgebraic set is semialgebraically path
connected~\cite[Theorems~5.22 and~5.23]{basupollackroy2006book}, so
connected and path components coincide. No step uses a bound on the
feasible set.
\end{proof}
